% CHECKPOINT35: formalization of the fully audited read-only normalization argument.
% Direct whole-B/R normalization supersedes the longer optional s-frame argument.
\statement{A.118}{Proved}{A normalized GT-zero residual involution and its transported pentagon}
Let \(\alpha\) be the original negative involution of full
\(\widehat{\mathrm{GT}}\), with canonical \(u\), parameter \(\phi\),
actual even \(k\), and \(d=1+2k\). Retain \(R\) and the specified
\(\Psi_R\) of A.116. There is an actual normalized
\(\Gamma=(d,f)\in\widehat{\mathrm{GT}}_0\) whose full modular
extension induces \(\Psi_R\), such that
\[
 \begin{gathered}
 \Gamma(\delta)=\delta^d,\qquad f=s\Gamma(s),\qquad
 f\equiv u\pmod R,\qquad f\in F',\\
 \Gamma(x)=x^d,\qquad \Gamma(y)=f^{-1}y^df.
 \end{gathered}                                                   \tag{A.118.1}
\]
For \(\beta=\Gamma^{-1}\alpha\Gamma\), its actual normalized
parameter and its unique canonical norm word are
\[
 \begin{gathered}
 \phi_\beta=\Gamma^{-1}\bigl(f^{-1}\phi\alpha(f)\bigr),\qquad
 U'=u_\beta=\Gamma^{-1}(f^{-1}u),\\
 \phi_\beta\in R,\qquad U'\in R,\qquad
 d_\beta=1,\qquad \beta=c_B\text{ on }B/R.
 \end{gathered}                                                   \tag{A.118.2}
\]
Here \(\beta\) is a negative involution of
\(\widehat{\mathrm{GT}}_0\); membership in full
\(\widehat{\mathrm{GT}}\) is not asserted.
The original full pentagon is retained exactly by the transported
identity (A.118.8) below. An ordinary pentagon for \(U'\) is not
inferred.

\begin{proof}
\emph{1. Preserve the specified whole modular quotient action.}
A.116 constructs a full peripheral automorphism \(\Gamma\) of \(B\)
inducing \emph{the specified} \(\Psi_R\), with
\(\Gamma(\delta)=\delta^d\) literally. Its construction first
left-composes with the inverse of an actual multiplier-\(d\) Galois
automorphism, uses A.3/A.34 with \(R=\overline{[R,B]}\), and then
composes that Galois automorphism back. The positive cusp transporter
in the lifting step belongs to \(R\). The final inner correction
therefore does not alter the action on \(B/R\).
Thus no additional arbitrary normalization of this full lift is
being performed here.

The induced map on \(S_3\) is the identity, so
\(f=s\Gamma(s)\) lies in \(F\). By the exact quotient assignment,
\[
 \Gamma(s)R=vR=suR,\qquad fR=uR.                              \tag{A.118.3}
\]
Since \(R\subset F'\) and \(u\in F'\), this already proves
\(f\in F'\). No extra prescription of an \(s\)-conjugator, and no
centralizer uniqueness in the quotient, is required.
The equation \(\Gamma(s)^2=1\) gives
\(\iota(f)=f^{-1}\), where \(\iota=\operatorname{Ad}(s)|_F\).
The original marking \(x=\delta^{-2}\), \(y=sxs\) now gives
\(\Gamma(x)=x^d\) and \(\Gamma(y)=f^{-1}y^df\).

The restriction of \(\Gamma\) commutes with the geometric
\(S_3\) in \(\operatorname{Out}(F)\): for each \(b\in B\),
\(\Gamma(b)\) and \(b\) have the same \(S_3\)-image, so their
conjugation actions on \(F\) have the same outer class.
The three named pure peripheral classes are the geometric orbit
of the first cusp; hence the common power is \(d\) on all three.
The actual Harbater--Schneps identification in A.1 therefore puts
this normalized \((d,f)\) in \(\widehat{\mathrm{GT}}_0\).
This argument uses the actual modular extension and its \(S_3\)
map. A generic peripheral lift without the specified whole quotient
action would not by itself give (A.118.3).

Since \(\Gamma\) is a full automorphism and induces the quotient
automorphism \(\Psi_R\), it preserves \(R\) in both directions.
Also \(\alpha\Psi_R=\Psi_Rc_B\) gives
\(\beta=c_B\) on \(B/R\).

\medskip\noindent\emph{2. The exact canonical norm under conjugation.}
The normalized multiplication rule is
\(f_{\gamma\eta}=f_\gamma\gamma(f_\eta)\).
It gives the displayed formula for \(\phi_\beta\). Directly,
\[
 \begin{aligned}
 \phi_\beta\,\beta\bigl(\Gamma^{-1}(f^{-1}u)\bigr)
 &=\Gamma^{-1}\bigl(
 f^{-1}\phi\alpha(f)\alpha(f^{-1}u)\bigr)\\
 &=\Gamma^{-1}\bigl(f^{-1}\phi\alpha(u)\bigr)
 =\Gamma^{-1}(f^{-1}u),
 \end{aligned}                                                   \tag{A.118.4}
\]
using the original norm equation \(u=\phi\alpha(u)\).
This solution belongs to \(F'\). The established triviality of
\(F^\alpha\) transfers to \(F^\beta\) under the actual full
\(\Gamma\). Two norm solutions in \(F'\) differ by a fixed
element, so the solution in (A.118.4) is the unique actual canonical
word \(U'\), rather than a candidate selected by a finite quotient.
Equation (A.118.3) and \(\Gamma(R)=R\) show \(U'\in R\).
Finally \(\phi_\beta=U'\beta(U')^{-1}\) lies in \(R\).

\medskip\noindent\emph{3. The complete dihedral normalization.}
This step retains the full dihedral quotient, rather than independent
primewise signs. Put
\[
 D=F/\overline{\langle\!\langle x^2,y^2\rangle\!\rangle}
   =\widehat{\mathbb Z}(T)\rtimes C_2(a),\qquad
 x\mapsto a,\quad y\mapsto aT.
\]
Its derived subgroup is \(2\widehat{\mathbb Z}(T)\).
For any normalized peripheral automorphism \(\gamma=(\lambda,f_\gamma)\),
write \(q_D(f_\gamma)=T^s\), with \(s\in2\widehat{\mathbb Z}\).
Peripheral preservation of \(x,y\) makes the square-normal kernel
\(\gamma\)-stable. Since \(\lambda\) is odd,
\[
 \gamma_D(a)=a,\qquad
 \gamma_D(T)=T^{D_\gamma},\qquad D_\gamma=1+2s.
\]
The original third peripheral power condition says that
\(T^{-D_\gamma}\) is conjugate to \(T^{-\lambda}\) in this full
semidirect product. A rotation has just its two global conjugates,
so \(D_\gamma=\epsilon\lambda\) for a single
\(\epsilon\in\{1,-1\}\). Since \(D_\gamma\equiv1\pmod4\),
\[
 D_\gamma=\chi_4(\lambda)\lambda,                            \tag{A.118.5}
\]
where \(\chi_4(\lambda)\in\{1,-1\}\) is the sign of the odd
unit modulo four. The sign is not chosen separately at different
primes.

Apply this to \(\Gamma\). The actual \(d\equiv1\pmod4\) gives
\(D_\Gamma=d\), so \(q_D(f)=T^{(d-1)/2}=T^k=q_D(u)\).
The exact formula for \(U'\) consequently gives \(q_D(U')=1\).
The original A.45 marking defines
\(q_D(U')=T^{k_\beta}\), \(d_\beta=1+2k_\beta\).
Faithfulness of the full rotation parameter proves
\(k_\beta=0\), hence \(d_\beta=1\).

For comparison, the general normalized conjugation formula is
\[
 \begin{gathered}
 u_{\gamma\alpha\gamma^{-1}}=f_\gamma\gamma(u_\alpha),\\
 d_{\gamma\alpha\gamma^{-1}}
 =\chi_4(\lambda_\gamma)\lambda_\gamma d_\alpha.
 \end{gathered}                                                   \tag{A.118.6}
\]
The norm identity and uniqueness prove the first line exactly as
in (A.118.4). Its dihedral exponent is
\(s+D_\gamma k\), giving the second line. No automorphy of a
prospective canonical endomorphism is used. Substitution of
\(\gamma=\Gamma^{-1}\) gives the same \(d_\beta=1\).

\medskip\noindent\emph{4. Preserve the full transported pentagon.}
Use the actual standard pure generators from A.94.1,
\[
 (a,b,c,d_{\rm br},e,f_{\rm br})
   =(A_{12},A_{13},A_{23},A_{14},A_{24},A_{34}),
\]
where \(A_{ij}\) are the generators denoted \(x_{ij}\) in the
original four-strand pentagon convention. Thus the parameter
\(f\) and the scalar \(d\) are unambiguous. In their original order set
\[
 \begin{gathered}
 (a_1,b_1)=(a,ce),\\
 (a_2,b_2)=(bc,f_{\rm br}),\qquad
 (a_3,b_3)=(c,f_{\rm br}),\\
 (a_4,b_4)=(ab,ef_{\rm br}),\qquad
 (a_5,b_5)=(a,c).
 \end{gathered}
\]
For each pair put
\[
 f_j=f(a_j,b_j),\qquad
 U'_j=U'\bigl(a_j^d,f_j^{-1}b_j^df_j\bigr).                 \tag{A.118.7}
\]
These evaluations are well-defined continuous homomorphisms from
\(F\) into the original profinite pure braid group. The exact word
identity \(u=f\Gamma(U')\) gives
\(u(a_j,b_j)=f_jU'_j\) separately for each of the five pairs.
The established full pentagon for this original canonical \(u\)
(A.37, with its full-\(\widehat{\mathrm{GT}}\) hypothesis) therefore
becomes precisely
\[
 \boxed{(f_1U'_1)(f_2U'_2)
       =(f_3U'_3)(f_4U'_4)(f_5U'_5).}                       \tag{A.118.8}
\]
No action of \(\Gamma\) on the whole four-strand group is assumed
in this substitution argument. The five occurrences of \(f_j\)
remain in their actual order; no ordinary pentagon for \(U'\) follows
by deleting them. In particular, membership of \(\Gamma\) or
\(\beta\) in full \(\widehat{\mathrm{GT}}\) has not been proved.
\end{proof}
\noindent\emph{Remaining obstruction.}
This reduces the original problem to an actual negative
\(\widehat{\mathrm{GT}}_0\)-involution which is complex conjugation
on \(B/R\), whose canonical word and parameter both lie in the
nontrivial residual \(R\), with actual dihedral parameter one and
the additional constraint (A.118.8). It neither kills that word nor
selects its full coherent frame class. Any already established
\(\widehat{\mathrm{GT}}_0\)-conjugacy invariant, including the original
marked CM obstruction at its proved scope, is retained under this
actual conjugation; no new value is inferred.
